% ---------------------------------------------------------------------------
\section{Considerações Finais}
% ---------------------------------------------------------------------------

O presente trabalho propõe atacar a deficiência inerente dos LLMs no raciocínio espacial contínuo transferindo os cálculos matemáticos para as \textit{game engines}. A arquitetura em desenvolvimento pauta-se no conceito de cooperação entre um componente estocástico-criativo (o modelo, como arquiteto) e um componente determinístico-rígido (o algoritmo, como engenheiro). 

À medida em que a prova de conceito (\textit{Addon} para Godot) avançar em testes empíricos de malhas navegáveis e consistência estrutural, será possível aferir a eventual redução da taxa de mapas inacessíveis, validando o arranjo teórico. Por fim, caso o método demonstre melhoras na solução problema investigado, projeta-se que a simbiose entre as descrições de \textit{layout} via LLM e algoritmos de construtivismo direcional fomentará metodologias de prototipagem hiper-ágeis para o mercado de jogos digitais, permitindo aos desenvolvedores concentrarem-se primordialmente na parte artística, enquanto delegam o trabalho estrutural bruto a modelos inteligentes aliados à lógica formal.
